\chapter{Metryzowalność przestrzeni topologicznych}
\label{ch:topog-metryzowalnosc}

\noindent\textbf{Cel rozdziału.}
Topologia przestrzeni metrycznej jest określona przez odległość,
lecz nie każda przestrzeń topologiczna posiada zgodną z topologią
metrykę. Ustalimy, jakie warunki są konieczne, a następnie
udowodnimy \emph{twierdzenie Urysohna o metryzowalności}:
przestrzeń regularna, spełniająca $T_1$ i drugi aksjomat
przeliczalności, jest metryzowalna. Nie ograniczymy się do
samego sformułowania: skonstruujemy rozdzielającą rodzinę funkcji
ciągłych i wyprowadzimy z niej metrykę. Dalej omówimy
metryzowalność podprzestrzeni i przeliczalnych iloczynów,
a także kontrprzykłady dotyczące iloczynów nieprzeliczalnych
i ilorazów.

\section{Metryki, pseudometryki i topologia wyznaczona przez odległość}
\label{sec:topog-metryki-definicje}

\begin{definicja}[Pseudometryka i metryka]
\label{def:topog-pseudometryka-metryka}
Niech $X$ będzie zbiorem. \emph{Pseudometryką} na $X$
nazywamy odwzorowanie $d\colon X\times X\to[0,\infty)$,
które dla wszystkich $x,y,z\in X$ spełnia
\begin{enumerate}[label=(\roman*)]
\item $d(x,x)=0$;
\item $d(x,y)=d(y,x)$;
\item $d(x,z)\le d(x,y)+d(y,z)$.
\end{enumerate}
Jeżeli ponadto $d(x,y)=0$ implikuje $x=y$, to $d$ jest
\emph{metryką}. Dla $x\in X$ i $r>0$ definiujemy kulę
$B_d(x,r)=\{y\in X:d(x,y)<r\}$.
Topologię generowaną przez wszystkie takie kule oznaczamy
$\tau_d$.
\end{definicja}

\begin{stwierdzenie}[Kule tworzą bazę]
\label{prop:topog-baza-kul}
Dla każdej pseudometryki $d$ na zbiorze $X$ rodzina
$\{B_d(x,r):x\in X,\ r>0\}$ jest bazą pewnej topologii
$\tau_d$. Dla każdego $x\in X$ rodzina
$\{B_d(x,1/n):n\in\mathbb N,\ n\ge1\}$ jest bazą otwartych
otoczeń $x$ w $\tau_d$.
\end{stwierdzenie}
\begin{proof}
Każdy punkt $x$ należy do $B_d(x,1)$, więc kule pokrywają
$X$; dla $X=\varnothing$ warunek jest pusty. Niech
$z\in B_d(x,r)\cap B_d(y,s)$. Liczba
$\delta=\min\{r-d(x,z),s-d(y,z)\}$ jest dodatnia.
Dla $w\in B_d(z,\delta)$ nierówność trójkąta daje
$d(x,w)\le d(x,z)+d(z,w)<r$ i analogicznie $d(y,w)<s$.
Stąd $B_d(z,\delta)$ mieści się w przecięciu kul, co jest
kryterium bazy topologii z rozdziału~01.

Jeśli $U\in\tau_d$ i $x\in U$, to istnieje $r>0$ z
$B_d(x,r)\subseteq U$: z bazowości mamy wprawdzie początkowo
$x\in B_d(z,s)\subseteq U$, lecz kula
$B_d(x,s-d(z,x))$ leży w $B_d(z,s)$. Wybierając
$n$ z $1/n<r$, otrzymujemy $B_d(x,1/n)\subseteq U$.
\end{proof}

\begin{definicja}[Metryzowalność]
\label{def:topog-metryzowalnosc}
Przestrzeń topologiczna $(X,\tau)$ jest \emph{metryzowalna},
jeżeli istnieje metryka $d$ na $X$, dla której $\tau=\tau_d$.
Taką metrykę nazywamy \emph{zgodną z topologią}.
Pojęcie nie odnosi się do konkretnego wzoru na odległość,
lecz do istnienia przynajmniej jednej zgodnej metryki.
\end{definicja}

\begin{przyklad}[Różnica między metryką a pseudometryką]
\label{ex:topog-pseudometryka-r2}
Na $\R^2$ funkcja
$p((x,y),(u,v))=|x-u|$ jest pseudometryką: nieujemność,
symetria i nierówność trójkąta wynikają z analogicznych
własności wartości bezwzględnej. Nie jest metryką, ponieważ
$p((0,0),(0,1))=0$. Każda kula $p$ jest pionowym pasem
$\{(u,v):|u-x|<r\}$, więc punkty o jednakowej pierwszej
współrzędnej należą dokładnie do tych samych zbiorów otwartych
$\tau_p$. W szczególności $\tau_p$ nie spełnia $T_0$.
Natomiast odległość euklidesowa na $\R^2$ jest metryką
i wyznacza zwykłą topologię euklidesową.
\end{przyklad}

\begin{stwierdzenie}[Warunki konieczne metryzowalności]
\label{prop:topog-metryka-t3-pierwszy}
Każda przestrzeń metryzowalna spełnia aksjomat $T_3$
oraz pierwszy aksjomat przeliczalności.
\end{stwierdzenie}
\begin{proof}
Niech $d$ będzie zgodną metryką. Pierwszy aksjomat
przeliczalności wynika z poprzedniego stwierdzenia o kulach.
Dla $x\ne y$ kula $B_d(x,d(x,y)/2)$ nie zawiera $y$;
ponieważ wszystkie kule są otwarte, singletony są domknięte:
$X\setminus\{y\}=\bigcup_{x\ne y} B_d(x,d(x,y)/2)$.
Otrzymujemy $T_1$.

Aby wykazać regularność, ustalmy domknięty $F\subseteq X$
i $x\notin F$. Zbiór $X\setminus F$ jest otwarty,
więc istnieje $r>0$ takie, że $B_d(x,r)\cap F=\varnothing$.
Pokażemy, że $\overline{B_d(x,r/3)}\subseteq B_d(x,2r/3)$.
Jeśli $y$ należy do lewego zbioru, to dla dowolnego
$\varepsilon>0$ istnieje $z\in B_d(x,r/3)$
z $d(y,z)<\varepsilon$; stąd
$d(x,y)<r/3+\varepsilon$.
Przechodząc z $\varepsilon\downarrow0$, otrzymujemy
$d(x,y)\le r/3<2r/3$. Zatem otwarte zbiory
$B_d(x,r/3)$ i $X\setminus\overline{B_d(x,r/3)}$
rozdzielają punkt $x$ i domknięty $F$. Jest to regularność,
a wraz z $T_1$ daje $T_3$.
\end{proof}

\begin{uwaga}[Ograniczenie warunków koniecznych]
Przestrzeń metryzowalna nie musi spełniać drugiego aksjomatu
przeliczalności. Niech $X$ będzie nieprzeliczalnym zbiorem,
a $d$ jego metryką dyskretną: $d(x,y)=0$ dla $x=y$
i $d(x,y)=1$ dla $x\ne y$. Dla każdej bazy $\mathcal A$
topologii $\tau_d$ i każdego $x\in X$, ponieważ
$\{x\}$ jest otwarty, istnieje $A\in\mathcal A$ z
$x\in A\subseteq\{x\}$. Wtedy $A=\{x\}$, zatem
każda baza zawiera wszystkie singletony i jest
nieprzeliczalna.
\end{uwaga}

\section{Metryki zgodne i podprzestrzenie}
\label{sec:topog-metryki-zgodne}

\begin{stwierdzenie}[Metryki o tych samych zbiorach otwartych]
\label{prop:topog-min-metryka}
Niech $d$ będzie metryką na zbiorze $X$ i niech
$\widehat d(x,y)=\min\{1,d(x,y)\}$.
Wtedy $\widehat d$ również jest metryką,
$\tau_{\widehat d}=\tau_d$ oraz $\widehat d\le1$.
\end{stwierdzenie}
\begin{proof}
Dodatniość i symetria są oczywiste. Dla $a,b\ge0$
zachodzi
$\min\{1,a+b\}\le\min\{1,a\}+\min\{1,b\}$:
jeżeli $a\ge1$ lub $b\ge1$, prawa strona jest co najmniej
$1$; w przeciwnym razie prawa strona wynosi $a+b$.
Stosując tę nierówność i nierówność trójkąta dla $d$,
uzyskujemy ją także dla $\widehat d$.
Dla każdego $0<r\le1$ i każdego $x$ zachodzi
$B_{\widehat d}(x,r)=B_d(x,r)$.
Kule o promieniach $1/n$, $n\ge2$, są bazami otoczeń
punktów w obu topologiach, więc topologie są równe.
\end{proof}

\begin{twierdzenie}[Dziedziczenie metryzowalności]
\label{thm:topog-metryzowalnosc-podprzestrzenie}
Niech $(X,\tau)$ będzie przestrzenią metryzowalną,
a $A\subseteq X$ jej podprzestrzenią. Wtedy topologia
podprzestrzeni na $A$ jest metryzowalna.
\end{twierdzenie}
\begin{proof}
Wybierzmy zgodną metrykę $d$ na $X$ i rozważmy jej
ograniczenie $d_A=d|_{A\times A}$. Wszystkie aksjomaty
metryki przechodzą na podzbiór. Dla $a\in A$ i $r>0$
mamy $B_{d_A}(a,r)=A\cap B_d(a,r)$.
Jeżeli $U\subseteq A$ jest otwarty w topologii podprzestrzeni,
ma postać $A\cap V$ dla otwartego $V\subseteq X$;
każdy jego punkt leży w odpowiedniej kuli $B_{d_A}(a,r)$.
Odwrotnie, każda taka kula jest względnie otwarta.
Zatem $\tau_{d_A}=\tau|_A$.
\end{proof}

\begin{stwierdzenie}[Przenoszenie metryki przez homeomorfizm]
\label{prop:topog-metryzowalnosc-homeo}
Jeżeli $h\colon X\to Y$ jest homeomorfizmem przestrzeni
topologicznych i $Y$ ma zgodną metrykę $d$, to
$d_h(x,x')=d(h(x),h(x'))$ jest metryką zgodną z topologią $X$.
\end{stwierdzenie}
\begin{proof}
Z różnowartościowości $h$ wynika dodatnia określoność,
a pozostałe aksjomaty dziedziczą się z $d$.
Dla $x\in X$ zachodzi
$B_{d_h}(x,r)=h^{-1}(B_d(h(x),r))$.
Ponieważ $h$ jest homeomorfizmem, przeciwobrazy kul
stanowią bazę topologii $X$.
\end{proof}

\section{Od regularności do normalności}
\label{sec:topog-t3-drugi-normalna}

Kluczowy etap konstrukcji metryki polega na otrzymaniu
\emph{przeliczalnie wielu} funkcji ciągłych rozdzielających
punkty i zbiory domknięte. Posłużymy się lematem Urysohna
udowodnionym w rozdziale~05. Wcześniej musimy uzasadnić
normalność, ponieważ regularność sama jej nie gwarantuje.

\begin{lemat}[Regularność i własność Lindelöfa dają normalność]
\label{lem:topog-t3-lindelof-normalna}
Niech $X$ będzie przestrzenią regularną spełniającą
własność Lindelöfa. Wtedy $X$ jest normalna. Jeżeli
ponadto $X$ spełnia $T_1$, to spełnia $T_4$.
\end{lemat}
\begin{proof}
Niech $F,G\subseteq X$ będą rozłącznymi zbiorami
domkniętymi. Jeżeli jeden z nich jest pusty, wystarczy
para $\varnothing,X$ ustawiona we właściwej kolejności.
Załóżmy więc, że oba są niepuste.

Dla każdego $x\in F$ regularność oraz kryterium zwężania
otoczeń z rozdziału~05 zapewniają otwarty $U_x\ni x$ taki,
że $\overline{U_x}\cap G=\varnothing$. Ponieważ zamknięta
podprzestrzeń przestrzeni Lindelöfa również spełnia
własność Lindelöfa (rozdział~09), możemy wybrać
ciąg $U_1,U_2,\ldots$ tych zbiorów pokrywający $F$.
Jeżeli podpokrycie jest skończone, powtarzamy jego elementy.
Analogicznie otrzymujemy otwarte $V_1,V_2,\ldots$ pokrywające
$G$ i spełniające $\overline{V_n}\cap F=\varnothing$.

Definiujemy zbiory otwarte
\[
U=\bigcup_{n\ge1}\left(U_n\setminus\bigcup_{k=1}^n
\overline{V_k}\right),\qquad
V=\bigcup_{n\ge1}\left(V_n\setminus\bigcup_{k=1}^n
\overline{U_k}\right).
\]
Każda różnica występująca po prawej jest otwarta,
bo odejmujemy skończoną sumę zbiorów domkniętych.
Dla $x\in F$ wybieramy $n$ z $x\in U_n$;
ponieważ $F\cap\overline{V_k}=\varnothing$ dla każdego $k$,
$x$ należy do $n$-tego składnika sumy definiującej $U$.
Stąd $F\subseteq U$. Symetrycznie $G\subseteq V$.

Pokażemy, że $U\cap V=\varnothing$. Gdyby $z\in U\cap V$,
istniałyby $n,m$ takie, że
$z\in U_n\setminus\bigcup_{k\le n}\overline{V_k}$
i $z\in V_m\setminus\bigcup_{k\le m}\overline{U_k}$.
Jeżeli $n\le m$, to $z\in U_n\subseteq\overline{U_n}$,
ale drugi warunek wyklucza przynależność do
$\overline{U_n}$. Jeżeli $m\le n$, otrzymujemy
analogiczną sprzeczność, korzystając z
$z\in V_m\subseteq\overline{V_m}$.
Tak więc $F$ i $G$ mają rozłączne otwarte otoczenia.
Otrzymaliśmy normalność, a wraz z $T_1$ — aksjomat $T_4$.
\end{proof}

\begin{wniosek}[Drugi aksjomat przeliczalności i $T_3$]
\label{cor:topog-drugi-t3-t4}
Każda przestrzeń spełniająca $T_3$ i drugi aksjomat
przeliczalności spełnia $T_4$.
\end{wniosek}
\begin{proof}
Z rozdziału~09 wynika, że drugi aksjomat przeliczalności
implikuje własność Lindelöfa. Aksjomat $T_3$ oznacza
regularność wraz z $T_1$. Stosujemy poprzedni lemat.
\end{proof}

\section{Twierdzenie Urysohna o metryzowalności}
\label{sec:topog-urysohn-metryzowalnosc}

\begin{twierdzenie}[Twierdzenie Urysohna o metryzowalności]
\label{thm:topog-urysohn-metryzowalnosc}
Niech $X$ będzie przestrzenią topologiczną spełniającą
aksjomat $T_3$ i drugi aksjomat przeliczalności.
Wówczas istnieje metryka $d$ na $X$ wyznaczająca
dokładnie jego wyjściową topologię.
\end{twierdzenie}
\begin{proof}
\textbf{Przypadki co najwyżej jednoelementowe.}
Jeśli $X=\varnothing$ lub $X$ ma jeden punkt, to
jedyna funkcja $d\colon X\times X\to[0,\infty)$
o wartości $0$ jest metryką, a wyznaczona przez nią
topologia jest jedyną możliwą topologią na $X$.
Dalej zakładamy, że $X$ ma co najmniej dwa punkty.

\textbf{Krok 1: przeliczalna rodzina par otwartych.}
Niech $\mathcal B$ będzie co najwyżej przeliczalną
bazą topologii $X$; można ją uznać za niepustą.
Rozważmy rodzinę
\[
\mathcal P=\{(B,C)\in\mathcal B\times\mathcal B:
\overline B\subseteq C\}.
\]
Rodzina $\mathcal P$ jest co najwyżej przeliczalna,
ponieważ jest podzbiorem iloczynu dwóch zbiorów
co najwyżej przeliczalnych.

Wykażemy, że dla każdego $x\in X$ i każdego
otwartego $O\ni x$ istnieje $(B,C)\in\mathcal P$ takie, że
\[
 x\in B\subseteq\overline B\subseteq C\subseteq O.
\]
Z bazowości $\mathcal B$ wybieramy $C\in\mathcal B$
z $x\in C\subseteq O$. Regularność i kryterium zwężania
otoczeń z rozdziału~05 dają otwarty $W$ z
$x\in W\subseteq\overline W\subseteq C$.
Ponownie z bazowości wybieramy $B\in\mathcal B$
z $x\in B\subseteq W$. Wtedy
$\overline B\subseteq\overline W\subseteq C$,
więc $(B,C)\in\mathcal P$. Ponieważ możemy
przyjąć $O=X\setminus\{y\}$ dla $x\ne y$ (aksjomat
$T_1$), rodzina $\mathcal P$ jest niepusta.

\textbf{Krok 2: funkcje oddzielające.}
Na mocy poprzedniego wniosku $X$ jest $T_4$.
Dla każdej pary $(B,C)\in\mathcal P$ zbiory
$\overline B$ i $X\setminus C$ są domknięte i rozłączne.
Lemat Urysohna z rozdziału~05 dostarcza zatem
funkcji ciągłej
$u_{B,C}\colon X\to[0,1]$ takiej, że
$u_{B,C}=0$ na $\overline B$ i $u_{B,C}=1$
na $X\setminus C$.
Wybieramy takie funkcje dla wszystkich par
z $\mathcal P$. Jest ich co najwyżej przeliczalnie wiele;
możemy ułożyć je w ciąg $(f_n)_{n\ge1}$,
powtarzając elementy, jeśli rodzina jest skończona.
Każda funkcja $f_n$ jest ciągła i przyjmuje wartości w $[0,1]$.

\textbf{Krok 3: konstrukcja i aksjomaty metryki.}
Dla $x,y\in X$ definiujemy
\[
 d(x,y)=\sum_{n=1}^{\infty}2^{-n}|f_n(x)-f_n(y)|.
\]
Szereg jest zbieżny, ponieważ jego wyrazy są
nieujemne i nie większe od $2^{-n}$, a
$\sum_{n\ge1}2^{-n}=1$. Otrzymujemy więc
$0\le d(x,y)\le1$, symetrię i $d(x,x)=0$.
Dla $x,y,z\in X$ nierówność trójkąta dla wartości
bezwzględnej, zastosowana do każdego składnika,
daje dla wszystkich $N$
\[
\sum_{n=1}^{N}2^{-n}|f_n(x)-f_n(z)|
\le\sum_{n=1}^{N}2^{-n}|f_n(x)-f_n(y)|
+\sum_{n=1}^{N}2^{-n}|f_n(y)-f_n(z)|.
\]
Po przejściu do granicy otrzymujemy
$d(x,z)\le d(x,y)+d(y,z)$.

Pozostała dodatnia określoność. Dla $x\ne y$
wybieramy w kroku 1 parę $(B,C)$ z
$x\in B$ i $C\subseteq X\setminus\{y\}$.
Odpowiednia funkcja $u_{B,C}=f_n$ spełnia
$f_n(x)=0$ i $f_n(y)=1$. Wobec tego
$d(x,y)\ge2^{-n}>0$. Zatem $d$ jest metryką.

\textbf{Krok 4: każda kula metryczna jest otwarta w $X$.}
Ustalmy $x\in X$ i $\varepsilon>0$. Wybierzmy $N$
tak duże, aby $\sum_{n>N}2^{-n}=2^{-N}<\varepsilon/2$.
Dla każdego $n\le N$, z ciągłości $f_n$ w punkcie $x$
zbiór
$W_n=\{y\in X:|f_n(y)-f_n(x)|<\varepsilon/2\}$
jest otwarty i zawiera $x$. Zbiór
$W=\bigcap_{n=1}^{N}W_n$ jest otwartym otoczeniem $x$.
Dla $y\in W$ mamy
\[
 d(x,y)\le\sum_{n=1}^{N}2^{-n}|f_n(x)-f_n(y)|
+\sum_{n>N}2^{-n}<\varepsilon/2+\varepsilon/2
=\varepsilon.
\]
A zatem $x\in W\subseteq B_d(x,\varepsilon)$.
Wykażmy teraz otwartość całej kuli, a nie tylko istnienie
otwartego otoczenia jej środka. Jeśli
$y\in B_d(x,\varepsilon)$, to
$\delta=\varepsilon-d(x,y)>0$. Powtarzając poprzednią
konstrukcję z punktem $y$ i promieniem $\delta$,
znajdujemy otwarty zbiór $W_y$ z
$y\in W_y\subseteq B_d(y,\delta)$.
Z nierówności trójkąta
$B_d(y,\delta)\subseteq B_d(x,\varepsilon)$.
Każdy punkt kuli ma więc zawarte w niej otwarte
otoczenie, czyli kula jest otwarta w początkowej
topologii. W konsekwencji $\tau_d\subseteq\tau$.

\pagebreak[4]
\textbf{Krok 5: każdy zbiór otwarty w $X$ jest otwarty
w metryce $d$.}
Niech $O\in\tau$ i $x\in O$. W kroku 1 znaleźliśmy
$(B,C)\in\mathcal P$ z $x\in B\subseteq C\subseteq O$.
Dla pewnego $n$ funkcja $f_n$ równa jest $0$ na $B$
i $1$ na $X\setminus C$. Jeśli $y\in X\setminus O$,
to $y\in X\setminus C$ oraz
$d(x,y)\ge2^{-n}|f_n(x)-f_n(y)|=2^{-n}$.
Wynika stąd
$B_d(x,2^{-n})\subseteq O$.
Ponieważ $x\in O$ był dowolny, $O\in\tau_d$,
czyli $\tau\subseteq\tau_d$.

Ostatecznie $\tau=\tau_d$, co kończy dowód.
\end{proof}

\begin{figure}[htbp]
\centering
\begin{tikzpicture}[>=Stealth, node distance=8mm,
 node/.style={draw,rounded corners=2pt,fill=blue!6,align=center,
 font=\small,text width=10.4cm,inner sep=6pt}]
\node[node] (a) {Baza $\mathcal B$ i regularność: dla $x\in O$
 dobieramy $B,C\in\mathcal B$ z $x\in B\subseteq\overline B\subseteq C\subseteq O$};
\node[node,below=of a] (b) {Normalność i lemat Urysohna:
 $f_{B,C}|_{\overline B}=0$, $f_{B,C}|_{X\setminus C}=1$};
\node[node,below=of b] (c) {$\displaystyle
 d(x,y)=\sum_{n\ge1}2^{-n}|f_n(x)-f_n(y)|$};
\node[node,below=of c] (d) {Metryka rozdziela punkty, a
 $\tau_d=\tau$};
\draw[->,thick,blue!65!black] (a)--(b);
\draw[->,thick,blue!65!black] (b)--(c);
\draw[->,thick,blue!65!black] (c)--(d);
\end{tikzpicture}
\caption{Etapy konstrukcji metryki w twierdzeniu Urysohna.
Przeliczalność bazy zapewnia, że do sumy wystarczy
ciąg funkcji $f_n$.}
\label{fig:topog-urysohn-metryka-schemat}
\end{figure}

\begin{wniosek}[Kryterium metryzowalności w klasie z drugim aksjomatem przeliczalności]
\label{cor:topog-drugi-metryzowalna-iff-t3}
Niech $X$ spełnia drugi aksjomat przeliczalności.
Wtedy $X$ jest metryzowalna wtedy i tylko wtedy, gdy
spełnia $T_3$.
\end{wniosek}
\begin{proof}
Jeżeli $X$ jest metryzowalna, to spełnia $T_3$
na mocy stwierdzenia \ref{prop:topog-metryka-t3-pierwszy}.
Odwrotna implikacja jest właśnie twierdzeniem
\ref{thm:topog-urysohn-metryzowalnosc}.
\end{proof}

\section{Kostka przeliczalna i przeliczalne iloczyny}
\label{sec:topog-metryzowalnosc-iloczyny}

\begin{twierdzenie}[Metryka na kostce przeliczalnej]
\label{thm:topog-metryka-kostka-przeliczalna}
Na $K=[0,1]^{\mathbb N}$ funkcja
\[
\rho(a,b)=\sum_{n=1}^{\infty}2^{-n}|a_n-b_n|,
\qquad a=(a_n),\ b=(b_n)\in K,
\]
jest metryką. Topologia $\tau_\rho$ jest dokładnie
topologią iloczynu na $K$.
\end{twierdzenie}
\begin{proof}
Zbieżność szeregu oraz aksjomaty metryki wynikają jak
w kroku 3 dowodu twierdzenia Urysohna, przyjmując
$f_n(a)=a_n$. Każde $a\ne b$ różni się na pewnej
współrzędnej $n$, więc $\rho(a,b)\ge2^{-n}|a_n-b_n|>0$.

Ustalmy $a\in K$ i $\varepsilon>0$.
Dobierzmy $N$ z $2^{-N}<\varepsilon/2$.
W topologii iloczynu jest otwarte otoczenie
\[
W=\{b\in K:|a_n-b_n|<\varepsilon/2
\text{ dla }1\le n\le N\}.
\]
Dla $b\in W$ suma pierwszych $N$ wyrazów definiujących
$\rho(a,b)$ jest mniejsza od $\varepsilon/2$,
a pozostałe wyrazy sumują się do co najwyżej $2^{-N}$.
Zatem $a\in W\subseteq B_\rho(a,\varepsilon)$.
Dla dowolnego $b\in B_\rho(a,\varepsilon)$
ponawiamy konstrukcję z promieniem
$\delta=\varepsilon-\rho(a,b)>0$ i środkiem $b$.
Otrzymujemy otwarty w topologii iloczynu zbiór
$W_b$ z $b\in W_b\subseteq B_\rho(b,\delta)$,
a nierówność trójkąta daje
$B_\rho(b,\delta)\subseteq B_\rho(a,\varepsilon)$.
Kula jest więc otwarta w topologii iloczynu,
czyli $\tau_\rho\subseteq\tau_{\mathrm{ilocz}}$.

Odwrotnie, niech $O$ będzie bazowym otwartym
otoczeniem $a$ w topologii iloczynu. Istnieje skończony
$F\subseteq\mathbb N$ i względnie otwarte
$O_n\subseteq[0,1]$, $n\in F$, takie, że
$a_n\in O_n$ i
$\{b\in K:b_n\in O_n\text{ dla }n\in F\}\subseteq O$.
Dla każdego $n\in F$ wybieramy $\delta_n>0$,
aby $\{t\in[0,1]:|t-a_n|<\delta_n\}\subseteq O_n$.
Jeśli $F\ne\varnothing$, przyjmijmy
$\delta=\min_{n\in F}2^{-n}\delta_n>0$.
Z $\rho(a,b)<\delta$ otrzymujemy, dla każdego $n\in F$,
$2^{-n}|a_n-b_n|\le\rho(a,b)<\delta\le2^{-n}\delta_n$.
Stąd $b_n\in O_n$, czyli $B_\rho(a,\delta)\subseteq O$.
Gdy $F=\varnothing$, możemy wziąć dowolne $\delta>0$,
ponieważ bazowy cylinder jest wtedy całym $K$.
Wynika stąd $\tau_{\mathrm{ilocz}}\subseteq\tau_\rho$.
\end{proof}

\begin{wniosek}[Reprezentacja w kostce Hilberta]
\label{cor:topog-urysohn-kostka-hilberta}
Każda przestrzeń $X$ spełniająca $T_3$ i drugi aksjomat
przeliczalności jest homeomorficzna z podprzestrzenią
kostki $[0,1]^{\mathbb N}$.
\end{wniosek}
\begin{proof}
Jeśli $X$ jest pusta, jej obraz pusty jest podprzestrzenią
kostki; dla przestrzeni jednoelementowej wystarcza odwzorowanie
stałe do ustalonego punktu kostki, które jest homeomorfizmem
na swój obraz.

Dla pozostałych przestrzeni wybierzmy funkcje $f_n$
z dowodu twierdzenia \ref{thm:topog-urysohn-metryzowalnosc}
i połóżmy
$\Phi(x)=(f_1(x),f_2(x),\ldots)$.
Wówczas dla wszystkich $x,y\in X$
\[
\rho(\Phi(x),\Phi(y))
=\sum_{n\ge1}2^{-n}|f_n(x)-f_n(y)|=d(x,y).
\]
Równość ta i metryczna charakteryzacja topologii
wykazana wyżej dowodzą, że $\Phi$ jest izometrią
z $(X,d)$ na $\Phi(X)$ z metryką ograniczoną z $\rho$.
Izometria jest różnowartościowa i zachowuje kule,
więc jest homeomorfizmem na swój obraz.
\end{proof}

\begin{twierdzenie}[Przeliczalne iloczyny przestrzeni metryzowalnych]
\label{thm:topog-przeliczalny-iloczyn-metryzowalny}
Niech $(X_n,\tau_n)$, $n\ge1$, będą przestrzeniami
metryzowalnymi. Wtedy iloczyn
$X=\prod_{n\ge1}X_n$, z topologią iloczynu,
jest metryzowalny. Konkretną zgodną metrykę można
skonstruować z dowolnych metryk $d_n$ zgodnych z $\tau_n$:
\[
D(x,y)=\sum_{n=1}^{\infty}
2^{-n}\min\{1,d_n(x_n,y_n)\}.
\]
\end{twierdzenie}
\begin{proof}
Jeżeli którykolwiek $X_n$ jest pusty, iloczyn jest pusty
i twierdzenie jest natychmiastowe. Zakładamy dalej,
że wszystkie czynniki są niepuste.
Niech $\widehat d_n=\min\{1,d_n\}$.
Poprzednie stwierdzenie
\ref{prop:topog-min-metryka} pokazuje, że każda
$\widehat d_n$ jest zgodną metryką ograniczoną przez $1$.

Szereg definiujący $D$ zbiega, ponieważ jego wyrazy
są ograniczone przez $2^{-n}$. Symetria, zerowanie
na przekątnej oraz nierówność trójkąta wynikają ze
zsumowania odpowiednich własności $\widehat d_n$.
Jeśli $x\ne y$, istnieje $k$ z $x_k\ne y_k$,
więc $D(x,y)\ge2^{-k}\widehat d_k(x_k,y_k)>0$.
Stąd $D$ jest metryką.

Ustalmy $x\in X$ i $\varepsilon>0$.
Wybierzmy $N$ tak, aby $2^{-N}<\varepsilon/2$.
Cylinder
\[
W=\{y\in X:\widehat d_n(x_n,y_n)<\varepsilon/2
\text{ dla }1\le n\le N\}
\]
jest otwartym otoczeniem $x$ w topologii iloczynu,
a dla $y\in W$ mamy $D(x,y)<\varepsilon$,
jak w dowodzie twierdzenia o kostce przeliczalnej.
W ten sposób każde $x\in X$ i każdy promień
$\varepsilon>0$ mają otwarte otoczenie $W$
zawarte w $B_D(x,\varepsilon)$. Aby stwierdzić,
że cała kula jest otwarta, dla dowolnego
$z\in B_D(x,\varepsilon)$ stosujemy ten sam argument
z promieniem $\varepsilon-D(x,z)$ i środkiem $z$.
Nierówność trójkąta zapewnia
$B_D(z,\varepsilon-D(x,z))\subseteq B_D(x,\varepsilon)$.
Stąd $\tau_D\subseteq\tau_{\mathrm{ilocz}}$.

Dla odwrotnej inkluzji rozważmy bazowy cylinder
$O=\prod_{n\ge1}O_n$ zawierający $x$,
gdzie $O_n$ jest otwarty w $X_n$, a $O_n=X_n$
poza pewnym skończonym zbiorem indeksów $F$.
Dla każdego $n\in F$ istnieje $r_n>0$ z
$B_{d_n}(x_n,r_n)\subseteq O_n$.
Przyjmijmy $\delta_n=\min\{1/2,r_n\}$.
Jeśli $D(x,y)<\delta=\min_{n\in F}2^{-n}\delta_n$,
to $\widehat d_n(x_n,y_n)<\delta_n<1$ i w konsekwencji
$d_n(x_n,y_n)=\widehat d_n(x_n,y_n)<r_n$ dla $n\in F$.
Tak więc $B_D(x,\delta)\subseteq O$.
Jeżeli $F=\varnothing$, cylinder jest całym $X$
i można przyjąć dowolny promień. Otrzymujemy
$\tau_{\mathrm{ilocz}}\subseteq\tau_D$, a więc
równość obu topologii.
\end{proof}

\begin{wniosek}[Zwarte przestrzenie Hausdorffa]
\label{cor:topog-zwarta-hausdorff-metryzowalna}
Niech $X$ będzie zwartą przestrzenią Hausdorffa.
Wtedy $X$ jest metryzowalna wtedy i tylko wtedy, gdy
spełnia drugi aksjomat przeliczalności.
\end{wniosek}
\begin{proof}
Jeżeli $X$ spełnia drugi aksjomat przeliczalności,
to na mocy rozdziału~07 zwarta przestrzeń Hausdorffa
jest normalna i $T_1$, a więc spełnia $T_3$.
Stosujemy twierdzenie Urysohna o metryzowalności.

Załóżmy teraz, że $X$ jest metryzowalna i wybierzmy
zgodną metrykę $d$. Dla każdego $n\ge1$ rodzina
$\{B_d(x,1/n):x\in X\}$ jest otwartym pokryciem $X$.
Ze zwartości możemy wybrać skończoną rodzinę
punktów $F_n\subseteq X$ taką, że
$X=\bigcup_{a\in F_n}B_d(a,1/n)$.
Jeśli $X$ jest pusta, te rodziny mogą być puste.
Zbiór $D=\bigcup_{n\ge1}F_n$ jest co najwyżej
przeliczalny. Jest gęsty w $X$: dla każdego
$x\in X$ i dowolnej kuli $B_d(x,r)$ wybierzmy
$n$ z $2/n<r$ oraz $a\in F_n$ z $d(x,a)<1/n<r$.
Zatem $D\cap B_d(x,r)\ne\varnothing$.
W rozdziale~09 udowodniliśmy, że ośrodkowość
przestrzeni metrycznej implikuje drugi aksjomat
przeliczalności. Stosujemy ten wynik do $(X,d)$.
\end{proof}

\section{Kontrprzykłady i granice twierdzeń}
\label{sec:topog-metryzowalnosc-kontrprzyklady}

\begin{przyklad}[Prosta Sorgenfreya: warunki konieczne nie wystarczają]
\label{ex:topog-sorgenfrey-niemetryzowalna}
Niech $S$ oznacza prostą Sorgenfreya z bazą
$\{[a,b):a<b\}$, omawianą w rozdziale~09.
Przestrzeń $S$ spełnia pierwszy aksjomat przeliczalności,
jest ośrodkowa i ma własność Lindelöfa, ale nie spełnia
drugiego aksjomatu przeliczalności; wszystkie te własności
zostały udowodnione w rozdziale~09.

Ponadto $S$ jest regularna i $T_1$: baza składa się
ze zbiorów zarazem otwartych i domkniętych.
Rzeczywiście dopełnienie $[a,b)$ jest sumą
$(-\infty,a)$ i $[b,\infty)$, a oba zbiory są otwarte
w topologii Sorgenfreya: pierwszy jest sumą bazowych
przedziałów $[x,a)$, $x<a$, a drugi sumą przedziałów
$[b,c)$ dla $c>b$ (w szczególności $[b,b+1)$,
$[b,b+2)$ itd. pokrywają $[b,\infty)$).
Dla dowolnego otwartego $O\ni x$ możemy dobrać
bazowy $B$ z $x\in B\subseteq O$, a
$\overline B=B\subseteq O$, więc kryterium
zwężania otoczeń z rozdziału~05 daje regularność.
Singletony są domknięte, ponieważ zwykła topologia
euklidesowa, spełniająca $T_1$, jest słabsza od
Sorgenfreya.

Gdyby $S$ była metryzowalna, jej ośrodkowość
implikowałaby drugi aksjomat przeliczalności,
na mocy twierdzenia o przestrzeniach metrycznych
z rozdziału~09. Prowadzi to do sprzeczności.
Zatem $S$ nie jest metryzowalna,
mimo spełniania $T_3$, pierwszego aksjomatu
przeliczalności i własności Lindelöfa.
\end{przyklad}

\begin{przyklad}[Nieprzeliczalna kostka]
\label{ex:topog-kostka-nieprzeliczalna}
Niech $I$ będzie zbiorem nieprzeliczalnym i
$K=[0,1]^I$ ma topologię iloczynu.
Na mocy twierdzenia Tychonoffa z rozdziału~07
$K$ jest zwarta; jest także Hausdorffa, gdyż każdy
czynnik $[0,1]$ jest Hausdorffa, a różniące się punkty
iloczynu można rozdzielić otoczeniami w jednej
współrzędnej.

Pokażemy, że $K$ nie spełnia nawet pierwszego
aksjomatu przeliczalności. Niech $0=(0)_{i\in I}$.
Przypuśćmy, że $(N_n)_{n\ge1}$ jest przeliczalną
bazą otoczeń $0$. Dla każdego $n$ wybierzmy
bazowy otwarty cylinder $W_n$ taki, że
$0\in W_n\subseteq N_n$. Istnieje skończony
zbiór współrzędnych $F_n\subseteq I$, poza którym
$W_n$ nie nakłada żadnych ograniczeń.
Suma $F=\bigcup_{n\ge1}F_n$ jest przeliczalna,
więc możemy wybrać $j\in I\setminus F$.

Zbiór $O_j=\{x\in K:x_j<1/2\}$ jest otwartym
otoczeniem $0$ w topologii iloczynu.
Istniałoby $n$ z $N_n\subseteq O_j$,
a wówczas $W_n\subseteq O_j$.
Ponieważ jednak $j\notin F_n$, punkt $z\in K$
zdefiniowany przez $z_j=1$ i $z_i=0$ dla $i\ne j$
należy do $W_n$, ale nie należy do $O_j$.
To sprzeczność.

Skoro każda przestrzeń metryzowalna spełnia
pierwszy aksjomat przeliczalności, $K$ nie jest
metryzowalna. W ten sposób otrzymujemy zwartą
przestrzeń Hausdorffa, która nie jest metryzowalna.
\end{przyklad}

\clearpage
\begin{przyklad}[Iloraz przestrzeni metryzowalnej nie musi być metryzowalny]
\label{ex:topog-iloraz-dwa-poczatki}
Weźmy rozłączną sumę dwóch kopii prostej rzeczywistej,
$Y=\R\times\{0,1\}$, z topologią sumy rozłącznej.
Jest metryzowalna: na $Y$ można zadać metrykę
\[
\delta((x,i),(y,j))=
\begin{cases}
\min\{1,|x-y|\},&i=j,\\
2,&i\ne j.
\end{cases}
\]
Nierówność trójkąta wynika z nierówności trójkąta
w każdej kopii oraz faktu, że każda odległość między
różnymi kopiami wynosi $2$, podczas gdy odległość
wewnątrz kopii nie przekracza $1$.
Topologia tej metryki jest sumą rozłączną: kule
o promieniu mniejszym od $1$ zawierają punkty tylko
jednej kopii i wyznaczają w niej zwykłą topologię.

Utożsammy $(t,0)$ z $(t,1)$ dla każdego $t\ne0$,
ale pozostawmy rozróżnione początki $(0,0)$ i $(0,1)$.
Niech $q\colon Y\to Z$ będzie odwzorowaniem ilorazowym.
Oznaczmy $o_0=q(0,0)$ i $o_1=q(0,1)$; mamy $o_0\ne o_1$.
Jeśli $U,V\subseteq Z$ są otwartymi otoczeniami
odpowiednio $o_0$ i $o_1$, to przeciwobrazy
$q^{-1}(U)$ i $q^{-1}(V)$ są otwarte w $Y$,
zawierając odpowiednio $(0,0)$ i $(0,1)$.
Istnieje więc $\varepsilon>0$ takie, że
$(t,0)\in q^{-1}(U)$ i $(t,1)\in q^{-1}(V)$
dla każdego $0<t<\varepsilon$.
Ponieważ $q(t,0)=q(t,1)$ dla $t\ne0$, otrzymujemy
$U\cap V\ne\varnothing$. Zatem dwóch początków
nie można rozdzielić rozłącznymi zbiorami otwartymi:
$Z$ nie jest Hausdorffa, a więc nie jest metryzowalna.
\end{przyklad}

\section{Zestawienie implikacji i zakresu założeń}
\label{sec:topog-metryzowalnosc-zestawienie}

\begin{stwierdzenie}[Podsumowanie kryteriów]
\label{prop:topog-metryzowalnosc-zestawienie}
Dla każdej przestrzeni topologicznej $X$ zachodzą następujące
stwierdzenia.
\begin{enumerate}[label=(\roman*)]
\item Jeśli $X$ jest metryzowalna, to spełnia $T_3$
i pierwszy aksjomat przeliczalności.
\item Jeśli $X$ spełnia $T_3$ i drugi aksjomat przeliczalności,
to jest metryzowalna.
\item Każda podprzestrzeń metryzowalnej przestrzeni
jest metryzowalna; to samo dotyczy przeliczalnego
iloczynu przestrzeni metryzowalnych.
\item Jeżeli $X$ jest zwarta i Hausdorffa, to
metryzowalność jest równoważna spełnianiu
drugiego aksjomatu przeliczalności.
\end{enumerate}
Żadnego z warunków $T_3$, pierwszego aksjomatu
przeliczalności, ośrodkowości i własności Lindelöfa
nie można, nawet po połączeniu wszystkich czterech,
użyć jako zamiennika drugiego aksjomatu przeliczalności
w twierdzeniu Urysohna.
\end{stwierdzenie}
\begin{proof}
Punkty (i)--(iv) są dokładnie wynikami
\ref{prop:topog-metryka-t3-pierwszy},
\ref{thm:topog-urysohn-metryzowalnosc},
\ref{thm:topog-metryzowalnosc-podprzestrzenie},
\ref{thm:topog-przeliczalny-iloczyn-metryzowalny}
i \ref{cor:topog-zwarta-hausdorff-metryzowalna}.
Ostatnią uwagę uzasadnia prosta Sorgenfreya z przykładu
\ref{ex:topog-sorgenfrey-niemetryzowalna}.
\end{proof}

\begin{uwaga}[Granice zakresu rozdziału]
Rozważaliśmy metryzowalność topologiczną,
nie zaś zupełność konkretnej metryki. Ta sama topologia
może pochodzić od wielu metryk o różnych własnościach
metrycznych. Na przykład na $\R$ zarówno
$d(x,y)=|x-y|$, jak i
$d'(x,y)=|\arctan(x)-\arctan(y)|$
indukują topologię euklidesową. Dla $d'$ wynika to
z faktu, że $\arctan\colon\R\to(-\pi/2,\pi/2)$
jest homeomorfizmem: jest ciągła, ściśle rosnąca,
a funkcja odwrotna $\tan$ jest ciągła na podanym
przedziale. Metryka $d'$ jest więc metryką przeniesioną
przez homeomorfizm na podprzestrzeń prostej metrycznej,
zgodnie ze stwierdzeniem
\ref{prop:topog-metryzowalnosc-homeo}.
Nie jest jednak zupełna: ciąg $x_n=n$ jest ciągiem
Cauchy'ego w metryce $d'$ (bo $\arctan(n)$ jest zbieżny),
natomiast nie ma granicy w $\R$ w żadnej z tych
równoważnych topologii. Metryzowalności nie wolno więc
utożsamiać ze zupełną metryzowalnością.
\end{uwaga}
